% Figure: an A-frame of two legs pinned at the apex, feet on the ground, a
% load at the apex. (a) the whole frame; (b) one leg isolated, showing the
% apex-pin interaction force that is internal to (a) and external to (b).
\begin{figure}[htbp]
\centering
\begin{tikzpicture}[
  force/.style={draw=engblue, -{Stealth}, very thick},
  rxn/.style={draw=engred, -{Stealth}, very thick},
  intn/.style={draw=engorange, -{Stealth}, very thick},
  scale=1.0,
]
% ===== (a) whole frame =====
\begin{scope}
  \coordinate (Ap) at (1.5,2.8);   % apex
  \coordinate (Lf) at (0,0);       % left foot
  \coordinate (Rf) at (3.0,0);     % right foot
  \draw[draw=engblack, very thick] (-0.4,0) -- (3.4,0);
  \foreach \x in {-0.3,0.1,...,3.3} {\draw[draw=enggray, thin] (\x,0) -- (\x-0.2,-0.26);}
  \draw[draw=engblack, line width=2.2pt] (Lf) -- (Ap) -- (Rf);
  \draw[draw=engblack, fill=enggray!20, thick] (Ap) circle (0.11);
  % load at apex
  \draw[force] (Ap) -- ++(0,-1.6) node[anchor=west, font=\small, engblue] {$W$};
  % foot reactions
  \draw[rxn] (Lf) -- ++(0,1.4) node[anchor=south, font=\small, engred] {$N_L$};
  \draw[rxn] (Rf) -- ++(0,1.4) node[anchor=south, font=\small, engred] {$N_R$};
  \node[font=\scriptsize, anchor=north east] at (Lf) {$L$};
  \node[font=\scriptsize, anchor=north west] at (Rf) {$R$};
  \node[font=\scriptsize, anchor=south] at (Ap) {$P$};
  \node[font=\small, anchor=north] at (1.5,-0.5) {(a) whole frame};
\end{scope}
% ===== (b) one leg isolated =====
\begin{scope}[xshift=6.2cm]
  \coordinate (Ap2) at (1.5,2.8);
  \coordinate (Lf2) at (0,0);
  \draw[draw=engblack, very thick] (-0.4,0) -- (2.2,0);
  \foreach \x in {-0.3,0.1,...,2.1} {\draw[draw=enggray, thin] (\x,0) -- (\x-0.2,-0.26);}
  \draw[draw=engblack, line width=2.2pt] (Lf2) -- (Ap2);
  \draw[draw=engblack, fill=enggray!20, thick] (Ap2) circle (0.11);
  % foot reaction
  \draw[rxn] (Lf2) -- ++(0,1.4) node[anchor=south, font=\small, engred] {$N_L$};
  % pin interaction from the other leg (external to this body)
  \draw[intn] (Ap2) -- ++(1.5,-0.6) node[anchor=west, font=\small, engorange] {$\vec{F}_{RL}$};
  \node[font=\scriptsize, anchor=north east] at (Lf2) {$L$};
  \node[font=\scriptsize, anchor=south] at (Ap2) {$P$};
  \node[font=\small, anchor=north] at (1.0,-0.5) {(b) left leg alone};
\end{scope}
\end{tikzpicture}
\caption[An A-frame and the apex-pin interaction]{(a) The whole A-frame, two
legs pinned at the apex $P$ with feet at $L$ and $R$ and a load $W$ hung at
the apex. Drawn as one body, the frame sees only external forces: the load
and the two foot reactions. The apex pin is internal and never appears. (b)
The left leg isolated as its own body. The cut through the apex pin now
exposes the force $\vec{F}_{RL}$ that the right leg exerts on the left, an
interaction force that was internal in (a) and is external in (b). By
Newton's third law the left leg pushes back on the right with $-\vec{F}_{RL}$,
which appears in the right leg's own diagram. The whole-frame diagram finds
the foot reactions; the single-leg diagram finds the apex-pin force the frame
diagram hides. A multi-body assembly is solved by choosing which cuts to make
and accepting the interaction forces those cuts expose.}
\label{fig:vol03ch01:a-frame}
\end{figure}
